\documentclass[11pt,a4paper]{article}
\usepackage[italian]{babel}
\usepackage[margin=2.5cm]{geometry}
\usepackage{parskip}
\usepackage{amsmath, amssymb, amsthm}
\usepackage{booktabs}
\usepackage{array}
\usepackage{xcolor}
\usepackage{needspace}
\usepackage{enumitem}
\usepackage{listings}
\usepackage{tikz}
\usepackage[most]{tcolorbox}
\usepackage[hidelinks]{hyperref}
\usetikzlibrary{decorations.pathreplacing, shapes.geometric, arrows.meta, positioning}

% Niente vedove/orfane: i paragrafi non lasciano righe isolate a fine/inizio pagina
\widowpenalty=10000
\clubpenalty=10000

% Elenchi più compatti
\setlist{itemsep=2pt, parsep=0pt, topsep=4pt, partopsep=0pt}

% --- Riquadri con bordo nero, senza numero (si spezzano tra due pagine) ---
% Uso: \begin{definizione} ... \end{definizione}
%      \begin{teorema}[Nome] ... \end{teorema}  (il nome è facoltativo)
\tcbset{nota/.style={enhanced jigsaw, breakable, colback=white,
  colframe=black, boxrule=0.5pt, sharp corners}}
\NewTColorBox{definizione}{}{nota,
  before upper={\textbf{Definizione.}\ }}
\NewTColorBox{teorema}{o}{nota,
  before upper={\textbf{Teorema\IfValueT{#1}{ (#1)}.}\ }}
\NewTColorBox{proposizione}{o}{nota,
  before upper={\textbf{Proposizione\IfValueT{#1}{ (#1)}.}\ }}
\NewTColorBox{proprieta}{o}{nota,
  before upper={\textbf{Proprietà\IfValueT{#1}{ (#1)}.}\ }}
\NewTColorBox{assioma}{o}{nota, boxrule=1pt,
  before upper={\textbf{Assioma\IfValueT{#1}{ (#1)}.}\ }}

% --- Ambienti semplici (senza riquadro) ---
% Il titolo sta in cima al blocco, anche se il contenuto inizia con un riquadro o
% con codice; e non resta mai da solo in fondo a una pagina.
% Uso: \begin{esempio}[Titolo facoltativo] ... \end{esempio}
\newcommand{\newrunin}[2]{%
  \NewDocumentEnvironment{#1}{o}{%
    \par\Needspace{4\baselineskip}\addvspace{\medskipamount}\noindent
    \textbf{#2}\IfValueT{##1}{ (##1)}\textbf{.}\ \ignorespaces}%
  {\par\addvspace{\medskipamount}}}
\newrunin{esempio}{Esempio}
\newrunin{esercizio}{Esercizio}
\newrunin{osservazione}{Osservazione}

% V e F nelle tavole di verità
\newcommand{\V}{\textsf{V}}
\newcommand{\F}{\textsf{F}}

% --- Codice C ---
% Uso: \begin{codice} ... \end{codice}      codice C numerato
%      \begin{comando} ... \end{comando}    comandi da terminale
%      \begin{risultato} ... \end{risultato} output di un programma
%      \lstinline|printf("ciao");|           codice dentro al testo
\lstdefinestyle{c}{
  language=C,
  basicstyle=\ttfamily\small,
  keywordstyle=\color{blue}\bfseries,
  commentstyle=\color{gray}\itshape,
  stringstyle=\color{red!70!black},
  numbers=left,
  numberstyle=\tiny\color{gray},
  numbersep=8pt,
  columns=flexible,
  keepspaces=true,
  breaklines=true,
  showstringspaces=false,
  tabsize=4,
  morekeywords={include, define},
  % lettere accentate nei commenti
  literate={à}{{\`a}}1 {è}{{\`e}}1 {é}{{\'e}}1 {ì}{{\`i}}1
           {ò}{{\`o}}1 {ù}{{\`u}}1 {È}{{\`E}}1 {À}{{\`A}}1
}
\lstset{style=c}
\newtcblisting{codice}{listing only, nota, left=6mm,
  listing options={style=c}}
\newtcblisting{comando}{listing only, nota,
  listing options={basicstyle=\ttfamily\small, numbers=none, language={},
    columns=flexible, keepspaces=true, breaklines=true}}
\newtcblisting{risultato}{listing only, enhanced jigsaw, breakable,
  colback=black!5, colframe=black, boxrule=0.5pt, sharp corners,
  listing options={basicstyle=\ttfamily\small, numbers=none, language={},
    columns=flexible, keepspaces=true, breaklines=true}}

% --- Diagrammi di flusso (simboli standard) ---
\tikzset{
  terminale/.style = {ellipse, draw, minimum width=2.5cm, minimum height=0.9cm},
  io/.style        = {trapezium, trapezium left angle=70, trapezium right angle=110,
                      draw, minimum width=2.5cm, minimum height=0.9cm},
  processo/.style  = {rectangle, draw, minimum width=2.5cm, minimum height=0.9cm},
  decisione/.style = {diamond, draw, aspect=2, minimum width=2.5cm},
  freccia/.style   = {thick, -{Stealth}}
}

\title{Sommatorie, ricorsione e polinomi}
\author{Marco Cerbella}
\date{Lezione 5 -- 5 ottobre 2026}

\begin{document}
\maketitle

\section{Sommatorie}

Sia $f \colon \mathbb{N} \to \mathbb{R}$ una successione, $f(n) = a_n$. Ad
esempio, per $a_n = n$ si ha $a_1 = 1$, $a_2 = 2$, \dots, $a_{100} = 100$.

Sommare i primi due termini significa calcolare $a_1 + a_2 = 1 + 2$; i primi
cinque, $a_1 + a_2 + a_3 + a_4 + a_5 = 1 + 2 + 3 + 4 + 5$. Per scrivere somme
con molti termini si usa il simbolo di \emph{sommatoria}:
\[
\sum_{i=1}^{n} f(i) = f(1) + f(2) + f(3) + \dots + f(n).
\]

\begin{esempio}
La somma dei primi $100$ numeri è
\[
1 + 2 + \dots + 100 = \sum_{i=1}^{100} a_i = \sum_{i=1}^{100} i .
\]
\end{esempio}

\begin{definizione}
Sia $a_n \colon \mathbb{N} \to \mathbb{R}$ una successione. Le si associa
un'altra successione
\[
s_k = a_1 + a_2 + \dots + a_k = \sum_{i=1}^{k} a_i .
\]
La successione $\{s_k\}$ si dice \emph{serie associata} alla successione
$\{a_n\}$.
\end{definizione}

\subsection{Proprietà della sommatoria}

\begin{proposizione}
Siano $\{a_n\}$, $\{b_n\}$ due successioni in $\mathbb{R}$ e $\bar n$ l'indice
iniziale della somma. Si hanno le seguenti proprietà:
\begin{enumerate}
    \item $\displaystyle \sum_{i=\bar n}^{n} (\lambda a_i)
          = \lambda \left( \sum_{i=\bar n}^{n} a_i \right)$, \quad
          $\lambda \in \mathbb{R}$;
    \item $\displaystyle \sum_{i=\bar n}^{n} (a_i + b_i)
          = \left( \sum_{i=\bar n}^{n} a_i \right)
          + \left( \sum_{i=\bar n}^{n} b_i \right)$;
    \item $\displaystyle \sum_{i=\bar n}^{n} a_i
          = \sum_{i=\bar n}^{k} a_i + \sum_{i=k+1}^{n} a_i$, \quad
          con $\bar n \leq k < n$;
    \item (cambio d'indice)
          $\displaystyle \sum_{i=1}^{n} a_i = \sum_{j=1}^{n} a_j
          = \sum_{k=1}^{n} a_k = \sum_{k=0}^{n-1} a_{k+1}$.
\end{enumerate}
\end{proposizione}

\section{Ricorsione}

\begin{definizione}
Una successione $\{a_n\}_{n \in \mathbb{N}}$ di numeri reali è
\emph{ricorsiva} se è definita specificando:
\begin{enumerate}
    \item il valore dei primi $m$ termini $a_1, a_2, \dots, a_m$;
    \item una funzione $f \colon \mathbb{R}^m \to \mathbb{R}$ tale che
          \[
            a_i = f(a_{i-1}, a_{i-2}, \dots, a_{i-m}) \qquad \forall\, i > m .
          \]
\end{enumerate}
\end{definizione}

\begin{esempio}
\begin{enumerate}
    \item Sia $a_0 = 0$ e $a_n = 1 + a_{n-1}$ per $n \geq 1$. Allora
          $a_1 = 1 + a_0 = 1$, $a_2 = 1 + a_1 = 2$, $a_3 = 1 + a_2 = 3$, \dots,
          $a_n = n$ (formula chiusa).
    \item \emph{Successione di Fibonacci}: $a_1 = 1$, $a_2 = 1$,
          $a_n = a_{n-1} + a_{n-2}$. Si ottiene $1, 1, 2, 3, 5, 8, \dots$
          (ad esempio $a_3 = a_2 + a_1 = 1 + 1 = 2$, $a_4 = a_3 + a_2 = 3$,
          $a_5 = a_4 + a_3 = 5$, $a_6 = a_5 + a_4 = 8$).
    \item Data $\{a_n\}$, la sua serie associata $\{s_n\}$ si definisce per
          ricorsione:
          \[
            \begin{cases}
              s_0 = a_0\\
              s_n = s_{n-1} + a_n
            \end{cases}
          \]
    \item $a_1 = 5$ e $a_n = 8$ per ogni $n > 1$: tutti i termini dopo il primo
          valgono $8$.
    \item \emph{Fattoriale}: $a_1 = 1$, $a_n = n \cdot a_{n-1}$. Allora
          $a_2 = 2 \cdot a_1 = 2$, $a_3 = 3 \cdot a_2 = 3 \cdot 2 \cdot 1$,
          $a_4 = 4 \cdot 3 \cdot 2 \cdot 1$, \dots, e in generale
          $a_n = n \cdot (n-1) \cdot (n-2) \cdots 3 \cdot 2 \cdot 1$.
          Si scrive $a_n = n!$ (formula esplicita), con $1! = 1$ e per
          convenzione $0! = 1$.
\end{enumerate}
\end{esempio}

\begin{definizione}
Data una successione $\{a_n\}$, una \emph{formula esplicita} o
\emph{formula chiusa} per $\{a_n\}$ è una formula $p(n)$ che esprime
direttamente il valore della successione $a_n$ a partire da $n$.

Se $\{a_n\}$ è definita per ricorsione, una formula chiusa per $\{a_n\}$ è
un'espressione $p(n)$ che verifica le condizioni iniziali e verifica la
relazione ricorsiva.
\end{definizione}

\begin{esempio}
Sia
\[
  \begin{cases}
    a_0 = 2\\
    a_n = a_{n-1} + 5
  \end{cases}
\]
Si ottiene $a_1 = 2 + 5 = 7$, $a_2 = 12$, $a_3 = 17$, $a_4 = 22$. La formula
chiusa $a_n = 5n + 2$ è quella giusta? Verifichiamo.
\begin{enumerate}
    \item \emph{Caso base} ($n = 0$): la formula dà $0 \cdot 5 + 2 = 2$, che
          coincide con il dato iniziale $a_0 = 2$.
    \item \emph{Relazione ricorsiva} $a_n = a_{n-1} + 5$: sostituendo la
          formula con $n-1$ al posto di $n$,
          \[
            a_{n-1} + 5 = \bigl( (n-1) \cdot 5 + 2 \bigr) + 5
            = 5n - 5 + 2 + 5 = 5n + 2 = a_n .
          \]
\end{enumerate}
Sono verificate entrambe: la formula chiusa è verificata.
\end{esempio}

\begin{esercizio}
Verificare che per
$\begin{cases} a_1 = 1\\ a_n = n + a_{n-1} \end{cases}$
la formula chiusa è $a_n = \dfrac{n(n+1)}{2}$.
\end{esercizio}

\emph{Svolgimento (aggiunto).} Per $n = 1$ la formula dà $\frac{1 \cdot 2}{2} = 1 = a_1$.
Inoltre
$a_{n-1} + n = \frac{(n-1)n}{2} + n = \frac{n^2 - n + 2n}{2} = \frac{n(n+1)}{2}$.

\subsection{Progressioni aritmetiche}

\begin{definizione}
Una successione $a_n \colon \mathbb{N} \to \mathbb{R}$ è una
\emph{progressione aritmetica} se
\[
  a_n - a_{n-1} = d, \qquad d \in \mathbb{R}.
\]
Definita in maniera ricorsiva, $a_0 = a$ e $a_n = d + a_{n-1}$, ha formula
chiusa
\[
  a_n = a + n \cdot d .
\]
\end{definizione}

\begin{esercizio}
Data $\begin{cases} a_0 = 8\\ a_n = a_{n-1} - 3 \end{cases}$, trovare la formula
chiusa e verificarla.
\end{esercizio}

\emph{Svolgimento (aggiunto).} È una progressione aritmetica con $a = 8$ e
$d = -3$, quindi $a_n = 8 - 3n$: per $n = 0$ vale $8$ e
$a_{n-1} - 3 = 8 - 3(n-1) - 3 = 8 - 3n = a_n$.

\section{Polinomi}

\begin{definizione}
Un \emph{polinomio nell'incognita $t$ a coefficienti in $\mathbb{R}$} è una
funzione $p \colon \mathbb{R} \to \mathbb{R}$ del tipo
\[
  p(t) = a_0 + a_1 t + a_2 t^2 + a_3 t^3 + \dots + a_n t^n
       = \sum_{i=0}^{n} a_i t^i ,
\]
dove $a_0, a_1, \dots, a_n \in \mathbb{R}$ sono i \emph{coefficienti} e $a_0$ è
il \emph{termine noto}.

Il più grande indice $k \in \mathbb{N}$ tale che $a_i = 0$ per $i > k$ e
$a_k \neq 0$ è il \emph{grado} del polinomio, $\operatorname{gr}(p) = k$. Se
$k$ non esiste, $\operatorname{gr}(p) = -\infty$.
\end{definizione}

\begin{esempio}
\begin{itemize}
    \item $o(t) = 0$ è il \emph{polinomio nullo}, con
          $\operatorname{gr}(o) = -\infty$.
    \item $q(t) = 3$ è un \emph{polinomio costante}
          ($a_0 = 3$, $a_1 = a_2 = \dots = 0$) e $\operatorname{gr}(q) = 0$.
    \item $p(t) = 5 - t + t^2$ ha $\operatorname{gr}(p) = 2$.
\end{itemize}
\end{esempio}

Si indica con $\mathbb{R}[t]$ l'insieme dei polinomi nell'incognita $t$ a
coefficienti reali, e con
\[
  \mathbb{R}_d[t] = \{ p \in \mathbb{R}[t] \mid \operatorname{gr}(p) \leq d \}
\]
quelli di grado al più $d$. Ad esempio $5 + t^2$ e $1 - 3t - t^2$ stanno in
$\mathbb{R}_2[t]$, il polinomio nullo $o$ sta in $\mathbb{R}_0[t]$, mentre
$5 + t^2 \notin \mathbb{R}_1[t]$.

\begin{osservazione}
Due polinomi $p = \sum_{i=0}^{n} a_i t^i$ e $q = \sum_{i=0}^{n} b_i t^i$ sono
uguali, $p = q$, \textbf{se e solo se} $a_i = b_i$ per ogni $i$. Ad esempio
$(x+2)^2 = x^2 + 4x + 4$.
\end{osservazione}

\subsection{Operazioni}

Dati $p = \sum_{i=0}^{n} a_i t^i$, $q = \sum_{i=0}^{n} b_i t^i$ e
$\lambda \in \mathbb{R}$, si definiscono:
\begin{itemize}
    \item il \emph{prodotto per uno scalare}:
          $(\lambda p) = \sum_{i=0}^{n} (\lambda a_i)\, t^i$;
    \item la \emph{somma}: $(p + q) = \sum_{i=0}^{n} (a_i + b_i)\, t^i$;
    \item il polinomio $t p = \sum_{i=0}^{n} a_i t^{i+1}
          = t a_0 + a_1 t^2 + \dots + a_n t^{n+1}
          = \sum_{i=1}^{n+1} a_{i-1} t^i$;
    \item il \emph{prodotto}:
          $p \cdot q = (a_0 + \dots + a_n t^n) \cdot q
          = a_0 q + a_1 t q + \dots + a_n t^n q$.
\end{itemize}

\begin{esempio}
Siano $p = 6 + 3t - t^2$ e $q = -6 + 5t$. Allora
\begin{align*}
p + q &= (6 - 6) + (3 + 5)\,t + (-1 + 0)\,t^2 = 8t - t^2,\\
p \cdot q &= (6 + 3t - t^2)(-6 + 5t)
  = 6(-6 + 5t) + 3t(-6 + 5t) - t^2(-6 + 5t)\\
  &= -36 + 30t - 18t + 15t^2 + 6t^2 - 5t^3
  = -36 + 12t + 21t^2 - 5t^3 .
\end{align*}
\end{esempio}

\begin{osservazione}
Per i gradi valgono
\[
  \operatorname{gr}(p + q) \leq \max\{\operatorname{gr}(p), \operatorname{gr}(q)\},
  \qquad
  \operatorname{gr}(p \cdot q) = \operatorname{gr}(p) + \operatorname{gr}(q).
\]
\end{osservazione}

\section{Divisione tra polinomi}

\begin{teorema}[Divisione con resto]
Siano $f, g \in \mathbb{R}[t]$ con $g \neq 0$. Allora esistono e sono unici due
polinomi $q, r \in \mathbb{R}[t]$ tali che
\[
  f = q \cdot g + r , \qquad \operatorname{gr}(r) < \operatorname{gr}(g).
\]
$q$ è il \emph{quoziente} e $r$ il \emph{resto}.
\end{teorema}

\begin{esempio}
Siano $f(t) = t^4 - t^3 + t - 1$ e $g(t) = t^2 - t$. Si ha
$t^2 \cdot g = t^4 - t^3$, quindi
\[
  f - t^2 g = t - 1 ,
\]
che ha grado $1 < \operatorname{gr}(g) = 2$: la divisione si ferma. Dunque
$q = t^2$, $r = t - 1$ e
\[
  f(t) = t^2 \cdot g + (t - 1).
\]
\end{esempio}

\begin{esempio}
Siano $f(t) = t^4 - t^3 + t - 1$ e $g(t) = t + 1$, con
$\operatorname{gr}(g) = 1$. Ad ogni passo si sottrae al resto parziale il
multiplo di $g$ che ne elimina il termine di grado massimo:
\[
\begin{array}{rcl}
  t^4 - t^3 + 0t^2 + t - 1 \;-\; t^3 g &=& -2t^3 + 0t^2 + t - 1 \\
  -2t^3 + 0t^2 + t - 1 \;-\; (-2t^2)\, g &=& 2t^2 + t - 1 \\
  2t^2 + t - 1 \;-\; (2t)\, g &=& -t - 1 \\
  -t - 1 \;-\; (-1)\, g &=& 0
\end{array}
\]
Il quoziente è $q = t^3 - 2t^2 + 2t - 1$ e il resto è $r = 0$
($\operatorname{gr}(0) = -\infty < \operatorname{gr}(g)$). Quindi
\[
  f(t) = (t^3 - 2t^2 + 2t - 1)\, g + 0 .
\]
\end{esempio}

\begin{definizione}
Siano $f, g \in \mathbb{R}[t]$ con $g \neq 0$. Si dice che $g$ \emph{divide}
$f$, e si scrive $g \mid f$, se esiste $q \in \mathbb{R}[t]$ tale che
$f = q \cdot g$ (cioè se il resto della divisione è nullo).
\end{definizione}

\end{document}
